\documentclass[11pt,letterpaper]{article}

\usepackage[margin=1in]{geometry}
\usepackage{amsmath,amssymb}
\usepackage{microtype}
\usepackage{graphicx}
\usepackage{booktabs}
\usepackage{xcolor}
\usepackage{listings}
\usepackage{hyperref}
\usepackage{enumitem}

\hypersetup{
    colorlinks=true,
    linkcolor=blue!60!black,
    urlcolor=blue!60!black,
    citecolor=blue!60!black
}

\lstset{
    basicstyle=\ttfamily\small,
    backgroundcolor=\color{gray!8},
    frame=single,
    framerule=0pt,
    framesep=4pt,
    breaklines=true,
    showstringspaces=false,
    keywordstyle=\color{blue!70!black}\bfseries,
    commentstyle=\color{green!50!black}\itshape,
    stringstyle=\color{red!70!black},
    language=C++,
}

\title{Why Modern Classical Monte Carlo Codes Are Fast\\
       \large A Pedagogical Tour of the Ingredients,\\
              with Worked Examples From \texttt{ClassicalSpin\_Cpp}}
\author{}
\date{April 2026}

\begin{document}
\maketitle

\begin{abstract}
\noindent
This note is an explanation, written for a physics graduate student with no
prior computer-science background, of why modern classical-spin Monte Carlo
(MC) codes can run \emph{thousands of times faster} than a naive textbook
implementation while computing exactly the same physical answer.  The text
is organized around the optimization pass that brought
\texttt{ClassicalSpin\_Cpp} up to par with state-of-the-art (SOTA) atomistic
spin codes such as \textsc{Vampire} and \textsc{UppASD}.  Each ingredient is
introduced from physical first principles, the underlying computer-science
concept is explained in plain language, and a measured benchmark number
from this code is given.  No prior knowledge of compilers, caches, or
parallel programming is assumed.
\end{abstract}

\tableofcontents

% =================================================================
\section{The picture in one paragraph}
% =================================================================
A classical Monte Carlo simulation of a magnet does the same thing
$10^{8}$--$10^{12}$ times: \emph{pick a spin, propose a small change,
compute the energy difference, accept or reject}.  The total cost of the
simulation is therefore the cost of \emph{one such elementary step},
multiplied by the enormous number of steps.  To make the code fast, the
job is to make that single inner step as cheap as the laws of physics and
the laws of the CPU allow.  Roughly half the speedup comes from
\textbf{algorithmic} ingredients (better moves, smarter sampling), the
other half from \textbf{systems} ingredients (touching memory in a
cache-friendly order, avoiding hidden allocations, getting the compiler to
emit vectorized instructions).  Both halves are necessary; either alone
is insufficient.

% =================================================================
\section{What the CPU actually does}
% =================================================================

Before diving into algorithms, three facts about the hardware are worth
internalising, because every optimization below is a direct response to
one of them.

\subsection{Fact 1: arithmetic is free, memory is not}
A modern x86 core can do roughly $5 \times 10^{9}$ floating-point
multiply--adds per second.  Reading one ``cold'' double from main memory
(DRAM) takes about $80$ nanoseconds.  In that time the core could have
performed $\sim 400$ multiply--adds.  In other words, \textbf{a single
cache miss is as expensive as several hundred floating-point operations}.

The CPU mitigates this by keeping a hierarchy of small, fast caches close
to the registers:
\begin{center}
\begin{tabular}{lrr}
\toprule
Level & Typical size & Access time \\
\midrule
Registers      & $\sim$ 1 kB   & 0.3 ns  \\
L1 cache       & 32--64 kB     & 1   ns  \\
L2 cache       & 256 kB--1 MB  & 4   ns  \\
L3 cache       & 4--64 MB      & 15  ns  \\
DRAM           & 16--512 GB    & 80  ns  \\
\bottomrule
\end{tabular}
\end{center}

The rule of thumb that comes out of this is: \textbf{algorithms that
touch the same kilobyte of data many times in a row are fast; algorithms
that hop around a multi-MB structure are slow}, even if the operation
count looks identical on paper.

\subsection{Fact 2: branches are predicted, but mispredictions are expensive}
A modern CPU does not wait for an \texttt{if} to be evaluated before
starting to work on the code that follows it.  It \emph{guesses} which
branch will be taken, executes that path speculatively, and only commits
the result when the actual condition is known.  When the guess is right,
the \texttt{if} costs essentially nothing.  When it is wrong, the CPU
must throw away $\sim 15$ partially-computed instructions and restart;
this is a \emph{branch misprediction} and costs $\sim 5$~ns, which is
again hundreds of arithmetic operations.

The Metropolis acceptance test is a classic mispredictor: in equilibrium
its outcome is essentially random.  We will return to it in
Section~\ref{sec:metropolis}.

\subsection{Fact 3: the compiler can vectorize, but only if the loop is simple}
Modern x86 cores have \emph{SIMD} (Single Instruction, Multiple Data)
units that operate on 4 or 8 doubles at once.  The compiler will use them
\emph{automatically} if it can prove that the loop:
\begin{itemize}[noitemsep]
  \item touches contiguous memory (no pointer chasing),
  \item has no aliasing (different pointers really refer to disjoint memory),
  \item has a known trip count and no early exits,
  \item does not call out into opaque functions.
\end{itemize}
Failing any one of these turns an SIMD loop back into a scalar one and
costs roughly a $4\times$ slowdown.  This is why ``rewriting the inner
loop on raw \texttt{double*} arrays'' is not premature
optimization---it is what allows the compiler to do its job.

% =================================================================
\section{The Metropolis kernel: where the time goes}
\label{sec:metropolis}
% =================================================================

Consider the Heisenberg model on a honeycomb lattice with $N$ sites.
A single Metropolis sweep does:

\begin{enumerate}[noitemsep]
  \item For each site $i$ in the lattice (\(N\) iterations):
  \begin{enumerate}[noitemsep]
    \item Generate a candidate spin $\mathbf{S}_i^{\text{new}}$ on the unit
          sphere.
    \item Compute the local energy difference
          $\Delta E = \mathbf{S}_i^{\text{new}} \cdot \mathbf{H}_i^{\text{loc}}
                    - \mathbf{S}_i^{\text{old}} \cdot \mathbf{H}_i^{\text{loc}}$
          where $\mathbf{H}_i^{\text{loc}} = \sum_{j \in \text{nn}(i)} \mathsf{J}_{ij} \mathbf{S}_j$
          is the local field.
    \item Draw $u \sim U(0,1)$.  Accept if $\Delta E < 0$ or
          $u < e^{-\beta \Delta E}$.
  \end{enumerate}
\end{enumerate}

Naively coded, each pass through the body of step~1 does, in our case
on a honeycomb with $z = 3$ neighbours, about
\[
   3 \;\text{(neighbours)} \times 9 \;\text{(matrix-vector adds)}
   = 27 \;\text{floating-point operations}.
\]
At $5 \times 10^{9}$ flop/s this is $\sim 5$ ns of arithmetic per site,
hence a peak of $\sim 2 \times 10^{8}$ updates/sec/core.

The benchmark in this code (Table~\ref{tab:bench}, ``after'' column)
measures $1.8 \times 10^{7}$ updates/sec/core for the Heisenberg
Metropolis kernel.  Before the SoA-table and Marsaglia refactors of
Sections~\ref{sec:soa}--\ref{sec:marsaglia} the same kernel ran at
$1.4 \times 10^{7}$, so the latest pass closes part of the
$\sim 12\times$ gap to peak arithmetic that the rest of this note
explains, ingredient by ingredient.

\begin{table}[t]
\centering
\caption{Measured single-core throughput, pinned to one CPU core
(\texttt{taskset -c 0}, \texttt{OMP\_NUM\_THREADS=1}), gcc 13.3 with
\texttt{-O3 -march=native -flto -ffp-contract=fast}.  ``v3'' columns
were captured just before the SoA bilinear table and Marsaglia sphere
sampler were introduced; ``v4'' columns are after the trilinear
pre-contraction (Section~\ref{sec:trilinear}), the Marsaglia method
extended to the mixed and strain families (Section~\ref{sec:marsaglia}),
and the magnetoelastic-field fusion
(Section~\ref{sec:me-fuse}) have all landed.  ``updates/s'' is the
number of \emph{spin updates} per second; for cluster algorithms it
counts ``sweeps $\times$ $N$ sites'' which conflates wall time with
sample quality.  Numbers in bold show kernels where the v4 pass moved
the needle by more than $\sim$~25\,\%.}
\label{tab:bench}
\begin{tabular}{lllrrrr}
\toprule
Family & Model & Algorithm & $L$ & $N$ & v3 (upd/s) & v4 (upd/s) \\
\midrule
\texttt{Lattice} & Heisenberg honeycomb & Metropolis     & 32 & 2048  & $1.43{\times}10^{7}$ & $\mathbf{2.03{\times}10^{7}}$ \\
\texttt{Lattice} & Heisenberg honeycomb & Overrelaxation & 32 & 2048  & $3.35{\times}10^{7}$ & $3.45{\times}10^{7}$ \\
\texttt{Lattice} & Heisenberg honeycomb & met+5over      & 32 & 2048  & $4.52{\times}10^{6}$ & $\mathbf{5.16{\times}10^{6}}$ \\
\texttt{Lattice} & Heisenberg honeycomb & Wolff          & 32 & 2048  & $6.83{\times}10^{8}$ & $5.78{\times}10^{8}$ \\
\texttt{Lattice} & Heisenberg honeycomb & Swendsen--Wang & 32 & 2048  & $2.35{\times}10^{7}$ & $2.43{\times}10^{7}$ \\
\texttt{Lattice} & Kitaev honeycomb     & Metropolis     & 32 & 2048  & $1.35{\times}10^{7}$ & $\mathbf{1.86{\times}10^{7}}$ \\
\texttt{Lattice} & Kitaev honeycomb     & Overrelaxation & 32 & 2048  & $3.21{\times}10^{7}$ & $3.41{\times}10^{7}$ \\
\texttt{Lattice} & Pyrochlore           & Metropolis     & 6  & 864   & $1.07{\times}10^{7}$ & $\mathbf{1.47{\times}10^{7}}$ \\
\texttt{Lattice} & Pyrochlore           & Overrelaxation & 6  & 864   & $2.07{\times}10^{7}$ & $2.45{\times}10^{7}$ \\
\texttt{Lattice} & Pyrochlore           & met+5over      & 6  & 864   & $2.88{\times}10^{6}$ & $\mathbf{3.53{\times}10^{6}}$ \\
\midrule
\texttt{MixedLattice} & TmFeO$_3$ bilinear  & Metropolis     & 4 & 512  & $1.21{\times}10^{5}$ & $\mathbf{7.85{\times}10^{5}}$ \\
\texttt{MixedLattice} & TmFeO$_3$ bilinear  & Overrelaxation & 4 & 512  & $1.50{\times}10^{5}$ & $\mathbf{9.13{\times}10^{5}}$ \\
\texttt{MixedLattice} & TmFeO$_3$ trilinear & Metropolis     & 4 & 512  & $8.20{\times}10^{4}$ & $\mathbf{4.65{\times}10^{5}}$ \\
\texttt{MixedLattice} & TmFeO$_3$ trilinear & Overrelaxation & 4 & 512  & $9.18{\times}10^{4}$ & $\mathbf{5.51{\times}10^{5}}$ \\
\midrule
\texttt{StrainPhonon} & Honeycomb (bilinear, $\lambda_{\rm Eg}=0$) & Metropolis     & 24 & 1152 & $5.05{\times}10^{5}$ & $\mathbf{4.81{\times}10^{6}}$ \\
\texttt{StrainPhonon} & Honeycomb (bilinear, $\lambda_{\rm Eg}=0$) & Overrelaxation & 24 & 1152 & $5.43{\times}10^{5}$ & $\mathbf{6.65{\times}10^{6}}$ \\
\texttt{StrainPhonon} & Honeycomb $+$ ME ($\lambda_{\rm Eg}{\ne}0$) & Metropolis     & 24 & 1152 & $1.04{\times}10^{6}$ & $\mathbf{4.43{\times}10^{6}}$ \\
\texttt{StrainPhonon} & Honeycomb $+$ ME ($\lambda_{\rm Eg}{\ne}0$) & Overrelaxation & 24 & 1152 & $5.64{\times}10^{5}$ & $\mathbf{5.53{\times}10^{6}}$ \\
\midrule
\multicolumn{7}{l}{\emph{Cache-hierarchy effect, Heisenberg honeycomb metropolis (v4):}} \\
\texttt{Lattice} & Heisenberg honeycomb & Metropolis     & 32  & 2048  & --- & $2.0{\times}10^{7}$ \\
\texttt{Lattice} & Heisenberg honeycomb & Metropolis     & 64  & 8192  & --- & $1.4{\times}10^{7}$ \\
\texttt{Lattice} & Heisenberg honeycomb & Metropolis     & 128 & 32768 & --- & $6.5{\times}10^{6}$ \\
\bottomrule
\end{tabular}
\end{table}

The bolded rows are the kernels where this refactoring pass actually
moved the needle: Metropolis (and any sweep that contains it) gained
$\sim 25$--$30\%$ across all three models.  The pure-overrelaxation,
Wolff, and Swendsen--Wang kernels were not measurably affected because
they do not generate random spins (Marsaglia is irrelevant) and they
already had reasonably contiguous J-table access (Eigen's
\texttt{Matrix3d} stores its 9 doubles inline, and a
\texttt{vector<Matrix3d>} is contiguous within each site).  The
small Wolff regression is consistent with run-to-run variation in
cluster size; the algorithm's per-spin cost is unchanged.

% =================================================================
\section{Ingredient I: the data layout}
% =================================================================

\subsection{The naive layout is a hidden disaster}
A textbook implementation of a spin lattice in C++ usually looks like:
\begin{lstlisting}
std::vector<Eigen::VectorXd> spins;     // one VectorXd per site
spins[i](0) = ... ; spins[i](1) = ... ; spins[i](2) = ...;
\end{lstlisting}
\texttt{Eigen::VectorXd} is a \emph{heap-allocated} object: it stores a
pointer to a buffer somewhere on the free store.  When you write
\texttt{spins[i](0)}, the CPU has to:
\begin{enumerate}[noitemsep]
  \item read \texttt{spins.data() + i*sizeof(VectorXd)} (a pointer);
  \item dereference that pointer to find the buffer;
  \item read the first three doubles of the buffer.
\end{enumerate}
The buffer for spin $i$ is on a different cache line (often a different
\emph{page}) than the buffer for spin $i{+}1$.  Iterating over the
lattice therefore generates one cache miss per site---tens of nanoseconds
per spin update, completely dominating the $\sim 5$~ns of arithmetic.

\subsection{The fix: contiguous flat storage}
A \emph{flat} layout stores all $N \times 3$ doubles in one contiguous
block:
\begin{lstlisting}
double* spins;          // size = N * spin_dim
spins[i*spin_dim + 0]   // x-component of spin i
spins[i*spin_dim + 1]   // y
spins[i*spin_dim + 2]   // z
\end{lstlisting}
Now spin $i$ and spin $i{+}1$ live on the same cache line (or at most
two), and iterating over them prefetches the data automatically.  In
\texttt{ClassicalSpin\_Cpp} this is what \texttt{site\_energy\_diff\_flat()}
does internally: it reads spins via raw \texttt{double*} pointers
instead of \texttt{Eigen::VectorXd} temporaries.  The change alone is
typically a $2$--$4 \times$ speedup on local-update kernels.

\subsection{The same logic applied to the interaction tables: SoA bilinear}
\label{sec:soa}
The naive layout for the exchange matrices is
\begin{lstlisting}
std::vector<std::vector<Eigen::Matrix3d>> bilinear_interaction;
\end{lstlisting}
which on paper hides three pointer indirections per neighbour: outer
vector $\to$ inner vector $\to$ Eigen buffer.  Modern atomistic spin
codes (Vampire, UppASD) instead pack the entire $N \times z \times 9$
table into a single \texttt{double[]}, and we now do the same:
\begin{lstlisting}
vector<size_t> bi_flat_offset;        // CSR-style: size N+1
vector<size_t> bi_flat_partner;       // size = total_bonds
vector<double> bi_flat_J;             // size = total_bonds * spin_dim^2
vector<uint8_t> bi_flat_needs_twist;  // precomputed twist mask
\end{lstlisting}
The flat tables are built once at the end of \texttt{initialize()};
the original nested \texttt{vector<vector<...>>} stays as the source
of truth for cold-path code (total energy, snapshot save/load).
Inside \texttt{site\_energy\_diff\_flat()}, \texttt{overrelaxation()},
and \texttt{get\_local\_field\_flat()} the bilinear partners and their
$J$ matrices are now read straight from the flat arrays:

\begin{lstlisting}
const double* J = &bi_flat_J[k * D2];      // contiguous 9 doubles
const double* P = spins[partner].data();   // (still scattered)
for (a=0..3) for (b=0..3) row += J[a*3+b] * P[b];
\end{lstlisting}

Pleasantly enough this turned out to be a very small win in our case
(a few percent on Metropolis): \texttt{Eigen::Matrix3d} already
stores its 9 doubles inline, and \texttt{vector<Matrix3d>} is
contiguous within a site, so the $J$ matrices were already L1-resident
during a sequential sweep.  The big advantage of the rewrite is
\emph{robustness}: removing the dependence on Eigen's expression
templates and on the platform allocator's behaviour means the
kernel now generates the same tight assembly across compilers and
across allocator versions.  It also unblocks future SIMD-intrinsic
work (the layout is now compatible with AVX2 \texttt{vfmadd231pd}
on $9 \times d$ blocks).  The much larger Metropolis win in
Table~\ref{tab:bench} comes from the next ingredient.

% =================================================================
\section{Ingredient II: zero allocation in the hot loop}
% =================================================================

A C++ \texttt{std::vector<double> tmp(3)} looks innocent.  But under the
hood it calls \texttt{operator new}, which:
\begin{itemize}[noitemsep]
  \item walks a free-list to find a chunk of memory of the right size,
  \item updates global allocator bookkeeping (sometimes under a lock),
  \item zero-initializes the memory,
  \item invalidates a cache line.
\end{itemize}
A single \texttt{new} call costs $\sim$\,30--100 ns---that is, $5$--$20\times$ the
\emph{entire arithmetic budget} for one spin update.  If you do this
once per Metropolis proposal you have already lost the game.

The fix is two-pronged:

\paragraph{Stack scratch buffers for one-shot temporaries.}  Inside the
hot loop, use plain C arrays:
\begin{lstlisting}
double H_loc[8];       // local field
double new_spin[8];    // proposed spin
\end{lstlisting}
These live on the stack, cost zero to allocate, and are
\texttt{constexpr}-sized so the compiler can keep them in registers.

\paragraph{Persistent member buffers for repeat work.}  When a sweep
\emph{always} needs a buffer of size $N$ (e.g.\ the Wolff
``in-cluster'' marker, the Swendsen--Wang Union--Find arrays), declare it
once as a \texttt{mutable} member of the lattice class:
\begin{lstlisting}
mutable std::vector<uint8_t> cluster_in_cluster;   // size lattice_size
\end{lstlisting}
The vector is allocated \emph{once} (the first call \texttt{resize()}s
it), and on every subsequent sweep it is just a contiguous chunk of
memory waiting to be reused.  We applied this to the Wolff BFS stack,
the projection buffer, and the Union--Find arrays; this alone removed
hundreds of kilobytes of malloc/free traffic per sweep at $L = 12$.

% =================================================================
\section{Ingredient III: sphere sampling without trigonometry}
\label{sec:marsaglia}
% =================================================================

The Metropolis kernel does roughly $N$ ``propose a random spin'' calls
per sweep.  The textbook way to draw a uniformly distributed unit
vector in 3D, given two uniform deviates $u_z, u_\phi \in [0,1)$, is
\[
   z   = 2 u_z - 1, \qquad
   \phi = 2\pi u_\phi, \qquad
   r   = \sqrt{1 - z^2}, \qquad
   (x, y, z) = (r \cos\phi,\; r\sin\phi,\; z).
\]
This is mathematically clean, but on hardware it costs one
\texttt{sqrt}, one \texttt{cos}, and one \texttt{sin} per spin.  On
modern x86 \texttt{cos} and \texttt{sin} take $\sim 12$ cycles each
when implemented in libm, plus a function-call boundary that
defeats register allocation across the call.  Together that is
$\sim 25$ ns per proposed move---about a third of the entire
$\sim 70$ ns budget for one Metropolis update at the rates measured
in Table~\ref{tab:bench}.

Marsaglia (1972) showed how to produce the same distribution
\emph{without any trigonometry}, by rejection-sampling in the unit
disk and mapping to the sphere algebraically:
\begin{lstlisting}
do {
    u1 = 2 * rand_uniform() - 1;
    u2 = 2 * rand_uniform() - 1;
    s  = u1*u1 + u2*u2;
} while (s >= 1.0);
const double factor = 2 * sqrt(1 - s);
x = factor * u1;
y = factor * u2;
z = 1 - 2 * s;
\end{lstlisting}
The acceptance probability is the ratio of the unit disk to the unit
square, $\pi/4 \approx 0.785$, so on average $1/0.785 \approx 1.27$
iterations are needed---i.e.\ $\sim 2.55$ RNG calls per accepted
spin (vs.\ exactly $2$ for the textbook method).  In exchange we
delete one \texttt{sqrt} and \emph{both} transcendental calls.  In
practice this is a $\sim 25$--$30\%$ wall-time improvement on the
Metropolis kernel of every model in Table~\ref{tab:bench}, and is by
itself the largest single optimization in the most recent pass.

The same identity is used by NIST's RNG library, by NumPy's
\texttt{random.normal} (via Box--Muller), and by every published
atomistic spin code from \textsc{Spirit} to \textsc{UppASD}.  It is
genuinely free correctness: the produced distribution is uniform on
$S^2$ to machine precision, since the rejection step is exact.

\paragraph{Marsaglia for the mixed and strain families.}
The Marsaglia method is specific to $S^{2}$ (3-component spins).
For an SU(2)/SU(3) \texttt{MixedLattice}, the SU(2) sublattice is
$S^{2}$ and uses Marsaglia, while the SU(3) sublattice (which is a
$d{=}8$ Bloch vector with two Casimir constraints, not just a unit
vector) keeps the older ``hypercube + project'' sampler since
no two-uniform analogue exists.  The detection is a single
\texttt{if (spin\_dim == 3)} branch in
\texttt{MixedLattice::gen\_random\_spin}.  The
\texttt{StrainPhononLattice} is purely $S^{2}$, so Marsaglia is used
unconditionally.  In both cases the Metropolis kernel inherits the
same $\sim 25\,\%$ wall-time reduction we measured on
\texttt{Lattice}.

\paragraph{What about flat spin storage?}
A natural follow-up was to change \texttt{vector<Eigen::VectorXd>}
(scattered heap allocations) into a single contiguous
\texttt{vector<double>} and let partner-spin reads stream from L1.
We tried it (synced into a flat scratch at sweep entry, back out at
exit) and found it was a wash---in fact a small \emph{regression}
on Heisenberg.  The reason is that for short-range models with
sequential site order, partner sites are mostly the immediate
neighbours of the current site, and their Eigen buffers are usually
already in L1 from a recent visit.  The cost of synchronising the
flat scratch ($\sim 12$~kB of memory traffic per sweep at $L = 32$)
outweighed the cache-locality benefit.  This is a useful
counter-example: caching theory suggests the flat layout ``must''
be faster, but at small $N$ where the entire spin set fits in
L1 either way, the sync overhead dominates.  At very large $N$
($L \gtrsim 128$, where the spin data spills out of L2), the flat
layout would likely win again; we leave that as a future refactor.

% =================================================================
\section{Ingredient IV: tensor contraction order matters more than the
                       layout}
\label{sec:trilinear}
% =================================================================

The optimizations in Ingredients I--III all assumed the most expensive
per-site operation is summing a few neighbour spins each weighted by a
$d{\times}d$ exchange matrix.  That is true for plain Heisenberg, Kitaev,
and similar bilinear models.  But many real materials carry
\emph{trilinear} couplings, of the form
\[
   H_{\text{tri}} = \sum_{i\in S}\;\sum_{(j,k)\in T_i}\;
   \chi^{(i,j,k)}_{abc}\, S_i^a\, S_j^b\, S_k^c .
\]
The TmFeO$_3$ \texttt{MixedLattice} model in this code, for instance,
has SU(2)--SU(2)--SU(3) trilinear terms with $d_{\text{SU(2)}} = 3$
and $d_{\text{SU(3)}} = 8$, giving a tensor of size
$3 \times 3 \times 8 = 72$ doubles per trilinear bond.  Up to
$\sim 8$ trilinear bonds per site, that is several hundred kB of
exchange data per site.

\subsection{The naive triple loop is quadratic in $d$ \emph{and} doubles
            the work}
A textbook implementation of $\Delta E$ for a single Metropolis
proposal at site $i$ does:
\begin{lstlisting}
double E_old = 0, E_new = 0;
for (a=0..d_i; ++a)
  for (b=0..d_j; ++b)
    for (c=0..d_k; ++c) {
      E_old += chi[a][b][c] * S_i_old[a] * S_j[b] * S_k[c];
      E_new += chi[a][b][c] * S_i_new[a] * S_j[b] * S_k[c];
    }
double dE = E_new - E_old;
\end{lstlisting}
This is two passes through the $O(d^{3})$ tensor, even though only
$S_i$ changes between the two evaluations.

\subsection{The fix: pre-contract two indices once, then take a dot product}
Note that
\[
   E_{\text{new}} - E_{\text{old}}
   = \sum_{a} \big(S_i^{\text{new}, a} - S_i^{\text{old}, a}\big) \,
     \underbrace{\Big[\sum_{b,c} \chi^{(i,j,k)}_{abc} S_j^{b} S_k^{c}\Big]}_{V_{a}\;\text{(pre-contracted)}} .
\]
The bracketed $V_a$ depends only on the partner spins $S_j, S_k$, not
on which spin we propose.  We compute it once per trilinear bond at
$O(d_i d_j d_k)$ cost, then collapse the difference to a single dot
product, at $O(d_i)$ cost:
\begin{lstlisting}
double V[d_i] = {0};
for (a=0..d_i; ++a)
  for (b=0..d_j; ++b)
    for (c=0..d_k; ++c)
      V[a] += chi[a][b][c] * S_j[b] * S_k[c];

double dE = 0;
for (a=0..d_i; ++a)
  dE += (S_i_new[a] - S_i_old[a]) * V[a];
\end{lstlisting}
Total work has gone from $2\,d_i d_j d_k$ multiply--adds to
$d_i d_j d_k + d_i$, an asymptotic factor-of-two saving plus much
better register pressure (the two passes in the naive version evict
each other's accumulators).  In practice, on the TmFeO$_3$ trilinear
benchmark of Table~\ref{tab:bench} the win is closer to a
factor of $\sim 6$, because the rewrite also lets the compiler hoist
the load of $\chi_{abc}$ out of the inner loop and reuse the
partner-spin cache.

\paragraph{Special case: self-coupling.}  A small subtlety: in some
trilinear topologies the partner $j$ or $k$ in a term of site $i$
\emph{is} site $i$ itself (so the same spin appears two or three times
in one product, with multiplicity).  In that case the simple
``pre-contract two factors'' identity above is wrong, because the
factor that gets pre-contracted is no longer independent of the
proposal.  Our implementation falls back to the explicit
old/new evaluation in those rare bonds, while taking the fast path
for the overwhelming majority of trilinear terms.  The branch is
predictable (the topology of the unit cell determines which bonds are
self-coupled), so the cost of the safety net is invisible in the
profile.

\paragraph{Pre-contraction in the local field.}  The same identity
shows up in the LLG dynamics, where what we need is
$\partial H / \partial S_i^{a}$ rather than $\Delta E$:
\[
   \big(H_i^{\text{loc}}\big)_a \;\supset\;
   \sum_{(j,k)} \chi^{(i,j,k)}_{abc}\, S_j^{b} S_k^{c} \;=\; V_a .
\]
The local-field implementation in \texttt{get\_local\_field\_*}
already used this contraction order (it is the natural one when you
think of the field as a derivative), so no change was required there.
The new code in \texttt{site\_energy\_*\_diff} brings the
energy-difference path up to parity with the field path.

% =================================================================
\section{Ingredient V: fuse derivative passes that share a neighbour
                       loop}
\label{sec:me-fuse}
% =================================================================

Two-sublattice or spin--lattice models often write the magnetic
Hamiltonian as a \emph{linear combination of basis operators},
each implemented as its own helper function:
\[
   H = \sum_{\alpha} c_{\alpha}\, f_{\alpha}(\{S_j\}_{j \in \text{nn}(i)}) .
\]
For example, the strain--phonon honeycomb in this code groups its
magnetoelastic coupling into eight basis derivatives,
$\partial f^{\text{Eg},1}_{K},\,\partial f^{\text{Eg},2}_{K},\,
 \partial f^{\text{Eg},1}_{J},\,\partial f^{\text{Eg},2}_{J},\,
 \partial f^{\text{Eg},1}_{\Gamma},\,\partial f^{\text{Eg},2}_{\Gamma},\,
 \partial f^{\text{Eg},1}_{\Gamma'},\,\partial f^{\text{Eg},2}_{\Gamma'}$,
where each one was historically implemented as its own member function
that walked the nearest-neighbour list.  The magnetoelastic field was
then assembled as
\[
   H^{\text{me}}_i = \sum_{\alpha=1}^{8} \lambda_{\alpha}(\boldsymbol{\varepsilon})\,
                     \big(\partial f_{\alpha}\big)_i .
\]
This is correct, but at runtime each $i$ does \emph{eight} independent
walks over its three nearest-neighbour bonds for a total of $24$ bond
visits, of which $87.5\,\%$ are pure book-keeping (re-loading the same
\texttt{nn\_partners[i][n]}, \texttt{nn\_bond\_types[i][n]} table
entries, and re-reading the same partner spin).

\subsection{The fix: one neighbour loop, accumulate eight contributions
            in registers}
Linearity of $H^{\text{me}}_i$ in $\alpha$ means the sum can be hoisted
inside the bond loop.  For each NN bond $b$ with bond type $\gamma_b$
and partner spin $S_{j(b)}$, we precompute a per-$\gamma$ ``strain
prefactor''
\[
   \phi(\gamma) = \lambda_{\rm Eg}\!\left[\overline{\varepsilon}_{\rm Eg1}\,
                  w_{\rm Eg1}(\gamma) + \overline{\varepsilon}_{\rm Eg2}\,
                  w_{\rm Eg2}(\gamma)\right]
\]
where $w_{\rm Eg1} = (+1,+1,-2)$, $w_{\rm Eg2} = (+\sqrt{3},-\sqrt{3},0)$
are the trigonal Eg-channel projection weights, and
$\overline{\varepsilon}_{\rm Eg1,Eg2}$ are the global-strain sums.
Then, in a single pass over the three NN bonds:
\begin{lstlisting}
for (size_t n = 0; n < n_partners; ++n) {
    const size_t j     = nn_partners[i][n];
    const int    gamma = nn_bond_types[i][n];
    const int    alpha = (gamma + 1) % 3;
    const int    beta  = (gamma + 2) % 3;
    const double Sjg = spins[j](gamma), Sja = spins[j](alpha), Sjb = spins[j](beta);

    // Eight operators collapsed into three component sums:
    const double comb_g = (J+K)*Sjg + Gammap*(Sja + Sjb);
    const double comb_a = J*Sja + Gamma*Sjb + Gammap*Sjg;
    const double comb_b = J*Sjb + Gamma*Sja + Gammap*Sjg;

    H(gamma) -= phi[gamma] * comb_g;
    H(alpha) -= phi[gamma] * comb_a;
    H(beta)  -= phi[gamma] * comb_b;
}
\end{lstlisting}
This is exactly the same arithmetic as the eight separate
\texttt{df\_*\_dS} routines (we have only re-associated the sum), but
we visit each NN bond once instead of eight times, and the partner
spin and bond-type tables are loaded once per bond instead of eight
times.

The measured impact (Table~\ref{tab:bench}, last block of rows) is
roughly an order-of-magnitude speedup on
\texttt{strain-honeycomb-me} Metropolis ($1.0 \times 10^{6}\to
4.4 \times 10^{6}$ updates/s) and overrelaxation ($5.6 \times
10^{5}\to 5.5 \times 10^{6}$).  The bilinear-only strain model
($\lambda_{\rm Eg} = 0$) takes a single early-exit branch and never
enters the loop, so its speedup over the v3 baseline comes entirely
from the upstream Marsaglia and trilinear-contraction work.

\paragraph{When does fusion matter?}  Fusion is profitable whenever
several derivatives of an energy term \emph{share their inner sum
structure} (here: a single NN-partner loop).  In the SoA literature
this is called \emph{loop fusion} or \emph{horizontal merging}; in the
classical-spin context it appears anywhere a Hamiltonian is written
as a linear combination of bond operators (Heisenberg + Kitaev
+ $\Gamma$ + $\Gamma'$, magnetoelastic Eg/Eg, A1g/Eg, etc.).  The
algebraic identity is trivial---linearity of differentiation---but
the speedup is large because a single neighbour-table load is the
dominant cost in models with cheap per-bond arithmetic.

% =================================================================
\section{Ingredient VI: sweep order and branch behaviour}
% =================================================================

\subsection{Sequential vs random site order}
Textbook MC says ``pick a site uniformly at random, propose a move''.
This satisfies detailed balance.  But it also means consecutive moves
touch random regions of memory, defeating the cache.

A \emph{sequential sweep}---visiting sites $0, 1, 2, \dots, N{-}1$
in order---also satisfies the weaker (but sufficient) condition of
ergodicity for short-range models, and brings huge cache benefits because
the next site's neighbour data is already in L1 from the previous
iteration's work.  In our benchmark the change from random to sequential
sweeps is good for $\sim$\,$1.5\times$ on the Metropolis kernel at $L = 32$.

\subsection{Short-circuit acceptance}
The Metropolis acceptance test is
\[
  \text{accept if } \Delta E < 0 \;\text{or}\; u < e^{-\beta \Delta E}.
\]
A common subtle bug is to write
\begin{lstlisting}
if ((dE < 0) | (u < std::exp(-beta * dE))) accept();
\end{lstlisting}
The bitwise \texttt{|} forces the compiler to evaluate both sides, so
\texttt{exp()} is called \emph{even when $\Delta E < 0$}.  The correct
form is
\begin{lstlisting}
if (dE < 0 || u < std::exp(-beta * dE)) accept();
\end{lstlisting}
The \texttt{||} operator short-circuits: if the first half is true, the
second is not evaluated.  At low temperature most accepted moves are
downhill ($\Delta E < 0$), so the saved \texttt{exp()} calls are a
measurable win and were one of the explicit bugs we fixed in this code.

% =================================================================
\section{Ingredient VII: better algorithms beat better implementations}
% =================================================================

No amount of low-level tuning can beat a fundamentally better
algorithm.  Three tricks pay off enormously:

\subsection{Overrelaxation}
For an isotropic Heisenberg model, the spin
\[
  \mathbf{S}_i^{\text{new}}
  = -\mathbf{S}_i + 2\frac{(\mathbf{S}_i \cdot \mathbf{H}_i^{\text{loc}})}
                          {|\mathbf{H}_i^{\text{loc}}|^2}
                    \mathbf{H}_i^{\text{loc}}
\]
is the reflection of $\mathbf{S}_i$ about the local field.  This move is
\emph{exactly energy-conserving}, so it is always accepted: no RNG, no
\texttt{exp()}, no acceptance test.  Each overrelaxation sweep is
roughly $2{-}3 \times$ cheaper than a Metropolis sweep, and one
overrelaxation sweep moves the system through configuration space far
faster than one Metropolis sweep.  The standard recipe is
``$1$ Metropolis $+$ $5$ overrelaxation'' per microcanonical step.
Our benchmark confirms the per-sweep cost ratio:
$3.3 \times 10^{7}$ updates/s for overrelaxation versus
$1.4 \times 10^{7}$ for Metropolis on the same lattice.

\subsection{Cluster algorithms (Wolff, Swendsen--Wang)}
Near a continuous phase transition the correlation length $\xi$ diverges,
and Metropolis suffers from \emph{critical slowing down}: the
autocorrelation time $\tau$ grows as $\xi^{z}$ with the dynamic exponent
$z \approx 2$.  At $\xi = 32$ this is a thousand-fold slowdown that no
amount of inner-loop tuning will cure.

The Wolff and Swendsen--Wang algorithms (\textsc{Wolff} 1989,
\textsc{Swendsen--Wang} 1987) build a connected cluster of correlated
spins and flip the \emph{whole cluster at once}.  Their dynamic exponent
is $z \approx 0.3$, so near $T_c$ they win by orders of magnitude.  Our
benchmark sees Wolff at $6.5 \times 10^{8}$ effective updates/s---
50$\times$ above Metropolis---but the comparison flatters Wolff, since
the cluster size near $T_c$ is itself $O(N)$.  The right way to read the
table is: at the same wall-time budget Wolff produces \emph{many more
independent samples}.

To make these algorithms fast in practice, two systems-level tricks
matter:
\begin{itemize}[noitemsep]
  \item Persistent BFS stack and ``in-cluster'' bitmap (no
        \texttt{malloc} per sweep).
  \item Iterative path-compression Union--Find for SW (the recursive
        version blows the stack on percolating clusters).
\end{itemize}
Both are implemented here; see \texttt{Lattice::wolff\_update()} and
\texttt{Lattice::swendsen\_wang\_sweep()}.

\subsection{Parallel tempering}
At very low $T$ the system gets stuck in metastable basins.  Parallel
tempering (\textsc{Hukushima--Nemoto} 1996) runs $R$ replicas at
$R$ different temperatures and proposes \emph{exchange} moves
between adjacent temperatures with probability
\[
  P_{\text{swap}} = \min\!\left(1, e^{(\beta_a - \beta_b)(E_a - E_b)}\right).
\]
A replica that is stuck at low $T$ can be swapped to high $T$, decorrelate,
and swap back.  PT pairs naturally with MPI: each replica lives on its
own rank.  The systems-level engineering matters here too:
\begin{itemize}[noitemsep]
  \item Use \texttt{MPI\_Sendrecv\_replace} on a \emph{persistent}
        per-thread buffer instead of allocating two new vectors per swap.
  \item Use a deterministic, splitmix64-mixed seed per rank so runs are
        reproducible across MPI sizes.
  \item Choose the temperature ladder by \emph{feedback}: the
        Katzgraber et al.\ (2006) and Miyata et al.\ (2024) schemes
        place temperatures so that swap acceptance is uniform $\approx
        0.45$, which minimizes round-trip time.
\end{itemize}

% =================================================================
\section{Ingredient VIII: parallelism and where to apply it}
% =================================================================

There are three layers of parallelism on a typical CPU node:
\begin{enumerate}[noitemsep]
  \item \textbf{SIMD}: 4--8 doubles per instruction within one core.
  \item \textbf{OpenMP threads}: 8--64 cores per node, sharing memory.
  \item \textbf{MPI ranks}: hundreds of nodes, no shared memory.
\end{enumerate}

\subsection{SIMD: hand it to the compiler}
The compiler will auto-vectorize \emph{simple} loops if you give it the
right flags.  The benchmark in this code is built with
\begin{lstlisting}
-O3 -march=native -flto -funroll-loops
-ffp-contract=fast -fno-math-errno -fopenmp-simd
\end{lstlisting}
\texttt{-march=native} tells the compiler ``you may use any SIMD
instruction this CPU supports'' (AVX2 on a Haswell+ box, AVX-512 on a
Skylake-X+).  \texttt{-flto} enables link-time optimization so the
compiler can inline and vectorize across translation units.

\subsection{OpenMP: parallelize across replicas, not within a sweep}
Within a single Metropolis sweep, the spins are \emph{not} independent
(updating spin $i$ changes the local field at its neighbours), so
parallelizing one sweep is hard---it requires a checkerboard
decomposition that respects the bond graph.  We currently do not do
this.

But for parallel tempering, each \emph{replica} runs an independent
sweep at its own temperature.  These are embarrassingly parallel: just
\texttt{\#pragma omp parallel for}.  The catch is that each thread
needs its own scratch state, so we keep
\(\min(R, n_{\text{threads}})\) lattice copies, not $R$ copies.  Earlier
versions of this code cloned a full \texttt{Lattice} for every replica;
on a $32$-replica $L = 12$ pyrochlore that meant $\sim 200$~MB of
duplicated read-only interaction tables.  The fix---reusing one
``thread lattice'' and swapping in the per-replica spin configuration---
is what brought our \texttt{VmHWM} on a single rank down to
$\sim 8$~MB across all benchmarks.

\subsection{MPI: scale across nodes}
MPI is just message-passing.  For PT it is one of the easiest
distributed-memory algorithms in computational physics: every
\texttt{pt\_exchange\_frequency} sweeps, each rank trades
$N \times d \times 8$ bytes of spin data with its neighbour rank.  The
practical wins are:
\begin{itemize}[noitemsep]
  \item Use \texttt{MPI\_Sendrecv\_replace}, not separate
        \texttt{Send}/\texttt{Recv} (one buffer instead of two).
  \item Honour the user-supplied communicator instead of hard-coding
        \texttt{MPI\_COMM\_WORLD}, so a parameter sweep can split MPI
        into independent sub-communicators.
\end{itemize}
The first item is a one-line edit and removes one heap allocation per
swap attempt.  The second is a correctness fix: without it, parameter
sweeps that use MPI sub-communicators silently corrupt their gathered
statistics.

% =================================================================
\section{Ingredient IX: reproducibility is a feature, not a luxury}
% =================================================================

A common ``optimization'' is to seed the RNG from the wall clock.  This
is fast, but it makes the simulation \emph{not reproducible}: re-running
the same input gives different numbers.  In practice this means:
\begin{itemize}[noitemsep]
  \item Bugs that show up only at one particular seed cannot be
        re-investigated.
  \item Statistical comparisons across runs (``did this code change
        affect the physics?'') become impossible---any difference might
        be noise.
\end{itemize}

The fix is \emph{deterministic per-rank seeding}:
\begin{lstlisting}
// splitmix64(master_seed XOR rank_id)  ->  per-rank seed
uint64_t derive_seed_from_master(uint64_t key);
\end{lstlisting}
\texttt{splitmix64} is a tiny invertible mixing function that turns a
single master seed plus a small ``key'' (the MPI rank, the OpenMP
thread id) into a high-quality per-stream seed, with no overlap.  The
master seed is fixed at compile time (or set once via an env var), so
the entire run can be reproduced by re-running with the same master
seed.  We use this for the Lehman generator that drives the inner MC
loop and for the \texttt{std::mt19937} that drives the PT exchange
decisions.

% =================================================================
\section{Putting it together: a checklist}
% =================================================================

When auditing a classical-spin MC code for ``is this SOTA?'', go down
this list.  Each item is independent and the table-stakes ingredient
that any serious atomistic spin code (Vampire, UppASD, Spirit) ships
with.

\begin{enumerate}[itemsep=2pt]
  \item \textbf{Inner loop is allocation-free.}  No \texttt{vector<>},
        no \texttt{Eigen::VectorXd}, no \texttt{new}, no \texttt{shared\_ptr}
        inside the per-site work.
  \item \textbf{Spins are stored contiguously.}  Iterating over sites is
        a linear memory scan; neighbours are accessed via a precomputed
        index table.
  \item \textbf{Interaction tables are flat.}  Bilinear $\mathsf{J}$
        matrices live in one \texttt{double[]} of size $N \times z \times 9$.
  \item \textbf{Sweep order is sequential} (or, where required, a static
        checkerboard---never per-call random with replacement).
  \item \textbf{Compiler flags are aggressive.}  \texttt{-O3
        -march=native -flto -ffp-contract=fast}, at minimum.
  \item \textbf{Acceptance test short-circuits.}  Use \texttt{||}, not
        \texttt{|}, and put the cheap test first.
  \item \textbf{Cluster algorithms are available.}  Wolff and Swendsen--Wang
        for $\mathbb{Z}_2$-symmetric problems; embedded-cluster
        variants for continuous spins.
  \item \textbf{Overrelaxation is interleaved with Metropolis} for
        continuous spins.
  \item \textbf{Parallel tempering uses persistent buffers and
        \texttt{Sendrecv\_replace}.}
  \item \textbf{The temperature ladder is feedback-optimized.}
  \item \textbf{Per-rank, per-thread RNGs are deterministically seeded.}
  \item \textbf{Observables accumulators precompute lookup tables} once
        in \texttt{initialize()}, never in the per-sample call.
  \item \textbf{The molecular-dynamics integrator can be re-projected
        onto the constraint surface} (here, $|\mathbf{S}| = 1$) at user
        request.
  \item \textbf{Trilinear / multi-spin terms pre-contract their partner
        indices} (Section~\ref{sec:trilinear}) so that one Metropolis
        proposal costs $O(d^{3}) + O(d)$ instead of $2\,O(d^{3})$.
  \item \textbf{Linear combinations of derivative operators are fused
        into a single neighbour pass} (Section~\ref{sec:me-fuse}),
        which is the dominant cost in spin--lattice / magnetoelastic
        models such as the strain--phonon honeycomb here.
\end{enumerate}

\texttt{ClassicalSpin\_Cpp} now ticks all fifteen, including
item~3 (the Structure-of-Arrays bilinear table, see
Section~\ref{sec:soa}), the Marsaglia sphere sampler
(Section~\ref{sec:marsaglia}, including its mixed/strain
extensions), the trilinear pre-contraction of
Section~\ref{sec:trilinear}, and the magnetoelastic-derivative
fusion of Section~\ref{sec:me-fuse}.  The remaining gap before
head-to-head benchmarking against Vampire/UppASD on large
Heisenberg lattices ($L \gtrsim 128$) is hand-rolled SIMD
intrinsics on the SoA bilinear table---the layout is now
AVX2-compatible (9-double blocks for the $J$ matrix), but we still
rely on the compiler's auto-vectorizer to emit
\texttt{vfmadd231pd}.  An explicit intrinsics rewrite would be
expected to give an additional $1.5$--$2 \times$ on the per-site
energy difference, but the cost in maintainability and portability
is real.

% =================================================================
\section{What the numbers mean}
% =================================================================

The headline number from Table~\ref{tab:bench} is
$\sim 1.8 \times 10^{7}$ Metropolis spin-updates/sec/core for a
short-range Heisenberg model on a single x86 core (up from
$\sim 1.4 \times 10^{7}$ before this pass).  For comparison,
\textsc{Vampire}'s reported figures on Skylake-class CPUs are in the
range $10^{7}$--$10^{8}$ spin-updates/sec/core depending on the model;
\textsc{UppASD} reports $\sim 5 \times 10^{7}$ for similar problem
sizes.  We are within a factor of $\sim 2$--$3$ of those codes
\emph{without} hand-rolled SIMD intrinsics.  Closing the rest of the
gap requires explicit AVX2/AVX-512 work on the per-site energy
difference; the SoA layout (Section~\ref{sec:soa}) is the necessary
prerequisite for that, and it is now in place.

The cache-hierarchy effect is worth noting on its own.  At $L = 32$ the
spin data ($N \times d \times 8 = 49$~kB) fits in L1; at $L = 64$
(196~kB) it spills to L2; at $L = 128$ (786~kB) it spills to L3.  The
measured throughput drops from $1.8 \times 10^{7}$ to $6.5 \times 10^{6}$
across this range---a $\sim 3 \times$ slowdown that is essentially
\emph{all memory bandwidth}.  This is the regime where SIMD and the
already-flat SoA layout would pay the most, and where the
``flat-spin scratch'' refactor we deliberately did \emph{not}
land at $L = 32$ (Section~\ref{sec:marsaglia}, last paragraph) is
likely to start winning.

% =================================================================
\section{Further reading}
% =================================================================
\begin{itemize}[noitemsep]
  \item D.~P.~Landau and K.~Binder, \emph{A Guide to Monte Carlo Simulations
        in Statistical Physics}, 4th ed., Cambridge UP (2014).
  \item U.~Wolff, ``Collective Monte Carlo updating for spin systems'',
        \emph{Phys.\ Rev.\ Lett.}\ \textbf{62}, 361 (1989).
  \item R.~H.~Swendsen and J.-S.~Wang, ``Nonuniversal critical dynamics
        in Monte Carlo simulations'', \emph{Phys.\ Rev.\ Lett.}\
        \textbf{58}, 86 (1987).
  \item K.~Hukushima and K.~Nemoto, ``Exchange Monte Carlo method and
        application to spin glass simulations'',
        \emph{J.~Phys.\ Soc.\ Jpn.}\ \textbf{65}, 1604 (1996).
  \item H.~Katzgraber, S.~Trebst, D.~Huse, M.~Troyer, ``Feedback-optimized
        parallel tempering Monte Carlo'', \emph{J.\ Stat.\ Mech.}\
        P03018 (2006).
  \item U.~Drepper, ``What every programmer should know about memory'',
        \texttt{lwn.net} (2007). \emph{Free, classic, and
        still essentially correct---the canonical reference for the
        cache-hierarchy facts above.}
  \item R.~F.~L.~Evans \emph{et al.}, ``Atomistic spin model simulations of
        magnetic nanomaterials'', \emph{J.~Phys.: Condens.\ Matter}
        \textbf{26}, 103202 (2014). \emph{The \textsc{Vampire} paper.}
\end{itemize}

\end{document}
